\section{Conclusion}

The Fibonacci trace map contains genuine symbolic, periodic-orbit, and
self-adjoint spectral structures, but their determinants are not
interchangeable.  A finite approximant discriminant is produced by a marked
chronological product over \(F_{k+2}\) sites.  A trace-map zeta coefficient is
a closed return after \(k\) renormalization steps.  The exact section/return
audit and the source-faithful ten-state boundary calculation expose this
semantic difference.

Two all-level results make the clock mismatch quantitative.  Under a uniform
local polynomial energy-degree bound, traces, marked boundary powers, and
order-\(k\) determinant coefficients have degree \(O(k)\), even if the finite
state dimension \(N_k\) grows arbitrarily.  The exact discriminant degree is
\(F_{k+2}\).  At the two rational finite-approximant section witnesses, the
discriminant values also grow so rapidly that their renormalization-clock
series has radius zero,
excluding literal realization by any fixed analytic germ at the origin.

The correct research consequence is a switch, not further local tuning.  A
future Fibonacci construction must visibly use physical time, growing-order
characteristic determinants, level-dependent data, a nonanalytic operator
mechanism, or a fully defined indirect divisor map.  Otherwise the next
exploration should move to a dynamical system whose arithmetic clock and
ordered cocycle are intrinsic from the start.
